\documentclass[a4paper,12pt]{article}
\usepackage[T2A]{fontenc}
\usepackage[utf8]{inputenc}
\usepackage[russian]{babel}
\usepackage{amsmath,amssymb,icomma}
\usepackage[margin=20mm,headheight=25pt,headsep=5mm,footskip=12mm]{geometry}
\usepackage{microtype,booktabs,tabularx,array,enumitem,fancyhdr,setspace,xcolor,tikz,titlesec}
\usetikzlibrary{arrows.meta,positioning,shapes.geometric}
\definecolor{accent}{HTML}{008080}
\definecolor{darkaccent}{HTML}{004D4D}
\definecolor{textgray}{HTML}{333333}
\color{textgray}
\onehalfspacing
\setlength{\parindent}{0pt}
\setlength{\parskip}{5pt}
\setlist{nosep,leftmargin=*}
\renewcommand{\arraystretch}{1.25}
\titleformat{\section}{\Large\normalfont\color{darkaccent}}{\thesection.}{.5em}{}[\vspace{4pt}\color{accent}\titlerule]
\titlespacing*{\section}{0pt}{12pt}{8pt}
\newcommand{\ex}[1]{\par\smallskip\textcolor{darkaccent}{\textbf{#1}}\par}
\newcommand{\kg}{\,\text{кг}}
\newcommand{\hours}{\,\text{ч}}
\pagestyle{fancy}
\fancyhf{}
\fancyhead[L]{\small Пропорции и проценты}
\fancyhead[R]{\small Надежда Клименко}
\fancyfoot[C]{\small\thepage}
\renewcommand{\headrule}{\hbox to\headwidth{\color{accent}\leaders\hrule height .5pt\hfill}}
\begin{document}
\begin{titlepage}
\centering
\vspace*{28mm}
{\Large\color{darkaccent} МАТЕМАТИКА\par}
\vspace{16mm}
{\Huge Чек-лист\par}
\vspace{10mm}
{\LARGE\color{accent} Прямая и обратная пропорциональность.\\[2mm] Задачи на проценты\par}
\vspace{14mm}
{\large Пропорции, метод стрелок и текстовые задачи\par}
\vfill
{\large Надежда Клименко\par}
\vspace{9mm}
25 ноября 2026 г.\par
\vspace{17mm}
{\small Лёвин Артём Александрович\\эксклюзивно для Надежды Клименко\par}
\vspace{19mm}
\begin{tikzpicture}
\draw[accent,line width=.7pt] (0,0) .. controls (2,0.75) and (4,-.75) .. (6,0) .. controls (8,.75) and (10,-.75) .. (12,0);
\end{tikzpicture}
\vspace{3mm}
\end{titlepage}
\setcounter{page}{1}
\section{Отношение и пропорция}
\textbf{Отношение} двух чисел $a$ и $b$ — их частное $a:b=\dfrac ab$, где $b\ne0$.

\textbf{Пропорция} — равенство двух отношений:
\[ a:b=c:d\quad\Longleftrightarrow\quad \frac ab=\frac cd,\qquad b\ne0,\ d\ne0. \]
Основное свойство: произведение крайних членов равно произведению средних:
\[\boxed{ad=bc}.\]
\begin{center}
\begin{tikzpicture}[>=Stealth]
\node (a) at (0,1.2) {$a$}; \node (b) at (0,0) {$b$};
\node (c) at (3,1.2) {$c$}; \node (d) at (3,0) {$d$};
\draw (-.35,.6)--(.35,.6);\draw (2.65,.6)--(3.35,.6);
\node at (1.5,.6) {$=$};
\draw[accent,thick] (a.east)--(d.west);\draw[accent,thick] (b.east)--(c.west);
\end{tikzpicture}
\end{center}
\ex{Пример 1. Можно ли составить пропорцию?}
Проверим отношения $3{,}06:0{,}9$ и $4{,}08:1{,}2$.
\[3{,}06\cdot1{,}2=3{,}672,\qquad 0{,}9\cdot4{,}08=3{,}672.\]
Произведения равны, значит, отношения образуют пропорцию.
\ex{Пример 2. Неизвестный член пропорции}
\[3\frac12:2\frac17=2\frac13:t,\qquad t\ne0.\]
Перемножаем крест-накрест, смешанные числа переводим в неправильные дроби:
\[\frac72t=\frac{15}{7}\cdot\frac73=5,\qquad t=5:\frac72=\frac{10}{7}=1\frac37.\]
При делении на дробь умножайте на обратную. Полученное значение $t$ отлично от нуля и удовлетворяет пропорции.
\clearpage
\section{Два вида пропорциональности}
Для положительных величин при фиксированных остальных условиях ($k>0$). Пусть $(a_1,b_1)$ и $(a_2,b_2)$ — две строки таблицы соответствующих значений.
\begin{center}
\begin{tabularx}{\linewidth}{@{}lXX@{}}
\toprule & Прямая & Обратная\\\midrule
Связь & $y=kx$ & $y=\dfrac{k}{x}$\\
Постоянно & $\dfrac yx=k$ & $xy=k$\\
Изменение $x$ в 2 раза & $y$ меняется в том же направлении в 2 раза & $y$ меняется в противоположном направлении в 2 раза\\
Стрелки & $\downarrow\quad\downarrow$ & $\downarrow\quad\uparrow$\\
Уравнение & $a_1b_2=a_2b_1$ & $a_1b_1=a_2b_2$\\\bottomrule
\end{tabularx}
\end{center}
\ex{Пример 3. Прямая пропорциональность}
Из $3{,}5$ кг семян получают $1{,}4$ кг масла. Сколько масла получится из $9{,}5$ кг семян при том же выходе?
\begin{center}
\begin{tikzpicture}[>=Stealth]
\node at (0,1.5) {Семена, кг}; \node at (4,1.5) {Масло, кг};
\node (a) at (0,.8) {$3{,}5$}; \node (b) at (4,.8) {$1{,}4$};
\node (c) at (0,-.2) {$9{,}5$}; \node (d) at (4,-.2) {$x$};
\draw[accent,->,thick] (-1,.9)--(-1,-.4);\draw[accent,->,thick] (5,.9)--(5,-.4);
\draw[accent] (a.east)--(d.west);\draw[accent] (b.west)--(c.east);
\end{tikzpicture}
\end{center}
\[3{,}5x=9{,}5\cdot1{,}4,\qquad x=\frac{9{,}5\cdot1{,}4}{3{,}5}=3{,}8\kg.\]
\ex{Пример 4. Обратная пропорциональность}
При скорости 60 км/ч путь занимает 4 ч. Сколько времени потребуется при скорости 80 км/ч на тот же путь?
\begin{center}
\begin{tikzpicture}[>=Stealth]
\node at (0,1.35) {Скорость, км/ч}; \node at (4,1.35) {Время, ч};
\node (a) at (0,.65) {$60$};\node (b) at (4,.65) {$4$};\node (c) at (0,-.15) {$80$};\node (d) at (4,-.15) {$x$};
\draw[accent,->,thick] (-1.45,.75)--(-1.45,-.35);\draw[accent,->,thick] (5,-.35)--(5,.75);
\draw[accent] (a.east)--(b.west);\draw[accent] (c.east)--(d.west);
\end{tikzpicture}
\end{center}
\[80x=60\cdot4,\qquad x=\frac{240}{80}=3\hours.\]
Проверка: $60\cdot4=80\cdot3=240$ км. При увеличении скорости время уменьшилось.
\clearpage
\section{Проценты: целое, часть и доля}
Один процент — одна сотая часть величины: $1\%=\dfrac{1}{100}$. Обозначим целое через $A>0$, часть через $B$, числовое значение процента через $p$.
\begin{center}
\begin{tabular}{@{}lcc@{}}
\toprule & Величина & Проценты\\\midrule
Целое & $A$ & $100$\\
Часть & $B$ & $p$\\\bottomrule
\end{tabular}
\end{center}
При фиксированном целом величина части прямо пропорциональна проценту. Основное уравнение: $100B=Ap$.
\begin{center}
\begin{tabularx}{\linewidth}{@{}lXl@{}}
\toprule Что найти & Правило & Формула\\\midrule
Часть & Умножить целое на долю & $B=\dfrac{Ap}{100}$\\[2pt]
Целое & Разделить часть на долю & $A=\dfrac{100B}{p}$, $p>0$\\[2pt]
Процент & Разделить часть на целое и умножить на 100 & $p=\dfrac{100B}{A}$\\\bottomrule
\end{tabularx}
\end{center}
\ex{Пример 5. Найти процент от целого}
Из 364 выпускников 91 продолжает музыкальное образование. Какой процент это составляет?
\[364x=91\cdot100,\qquad x=\frac{91\cdot100}{364}=25.\]
Ответ: $25\%$. Проверка: четверть от 364 равна 91.
\ex{Пример 6. Найти часть по проценту}
В $3{,}5$ кг технического золота содержится $41{,}5\%$ примесей. Найдём массу примесей:
\begin{center}
\begin{tabular}{@{}lcc@{}}
\toprule & Масса, кг & Проценты\\\midrule
Техническое золото & $3{,}5$ & $100$\\
Примеси & $x$ & $41{,}5$\\\bottomrule
\end{tabular}
\end{center}
\[100x=3{,}5\cdot41{,}5,\qquad x=1{,}4525\kg.\]
Проверка: масса примесей меньше общей массы. Единица ответа — килограммы.
\clearpage
\section{Как восстановить целое по его части}
\ex{Пример 7. Выход продукта}
В оливках содержится $64{,}8\%$ масла. Сколько килограммов оливок нужно для получения $40{,}5$ кг масла, если извлекается всё содержащееся масло?
\begin{center}
\begin{tabular}{@{}lcc@{}}
\toprule & Масса, кг & Проценты\\\midrule
Оливки & $x$ & $100$\\
Масло & $40{,}5$ & $64{,}8$\\\bottomrule
\end{tabular}
\end{center}
\[64{,}8x=40{,}5\cdot100,\qquad x=\frac{4050}{64{,}8}=62{,}5\kg.\]
Проверка: $62{,}5\cdot0{,}648=40{,}5$ кг.
\ex{Пример 8. Раствор и растворённое вещество}
Масса соли в растворе равна $3{,}6$ кг, массовая доля соли — $12\%$. Найдём массу всего раствора.
\begin{center}
\begin{tabular}{@{}lcc@{}}
\toprule & Масса, кг & Проценты\\\midrule
Раствор & $x$ & $100$\\
Соль & $3{,}6$ & $12$\\\bottomrule
\end{tabular}
\end{center}
\[12x=3{,}6\cdot100,\qquad x=30\kg.\]
\section{Когда процент может быть больше 100?}
В задачах на состав вещества или долю группы часть целого составляет от $0\%$ до $100\%$. При сравнении с выбранной базой результат может превышать $100\%$.

Например, если план составляет 40 изделий, а изготовлено 50, то выполнение плана равно $\dfrac{50}{40}\cdot100\%=125\%$. План перевыполнен на $25\%$.
\clearpage
\section{Типичные ошибки и самопроверка}
\begin{tabularx}{\linewidth}{@{}>{\raggedright\arraybackslash}p{.35\linewidth}X@{}}
\toprule Ошибка & Как исправить\\\midrule
Сравнение разных единиц & Под килограммами пишите килограммы, под процентами — проценты. При необходимости переведите величины в одинаковые единицы.\\
Неверный выбор 100\% & Определите целое или базу сравнения. Для раствора это масса всего раствора; для выхода продукта — масса исходного сырья.\\
Неправильное направление стрелок & Проверьте, что остаётся постоянным: отношение или произведение. Одного наблюдения «больше — больше» недостаточно.\\
Неправильное умножение & Для прямой зависимости — крест-накрест, для обратной — по строкам.\\
Ошибка при делении на дробь & Умножайте на обратную дробь; смешанные числа сначала преобразуйте.\\
Неверная единица ответа & Если $x$ стоит в столбце массы, ответ в килограммах; если в столбце процентов — в процентах.\\
\bottomrule
\end{tabularx}
\ex{Пример 9. Проверка базы сравнения}
Из 220 кг свежих абрикосов получили 55 кг кураги. За $100\%$ принимаем исходные 220 кг.
\[p=\frac{55\cdot100}{220}=25.\]
Выход кураги составляет $25\%$ от массы свежих абрикосов. Это именно отношение масс; по этим данным нельзя отдельно определить содержание сухого вещества в кураге.
\ex{Контроль перед записью ответа}
\begin{itemize}[label=$\square$]
\item Выбраны величины и одинаковые единицы измерения.
\item Установлено, какая зависимость действует и при каких условиях.
\item В задачах на проценты правильно определена база $100\%$.
\item Уравнение соответствует направлению стрелок.
\item Ответ проверен подстановкой и по смыслу задачи.
\end{itemize}
\clearpage
\section*{Алгоритм: решение задачи через пропорцию}
\vfill
\begin{center}
\begin{tikzpicture}[>=Stealth,every node/.style={font=\small},
act/.style={draw=accent,rounded corners=2pt,align=center,text width=10cm,inner sep=7pt},
term/.style={draw=accent,ellipse,align=center,inner sep=6pt},
dec/.style={draw=accent,diamond,aspect=2.6,align=center,inner sep=2pt},
branch/.style={draw=accent,rounded corners=2pt,align=center,text width=4.1cm,minimum height=2.1cm,inner sep=5pt},
arr/.style={->,thick,draw=darkaccent}]
\node[term] (start) at (0,0) {Прочитайте условие};
\node[act] (tab) at (0,-1.7) {Выделите две величины. Запишите данные в таблицу, приведите единицы к одному виду.};
\node[act] (base) at (0,-3.6) {Если задача на проценты, найдите базу $100\%$. Обозначьте неизвестное через $x$.};
\node[dec] (d) at (0,-5.65) {Отношение постоянно?};
\node[dec] (e) at (0,-8.1) {Произведение постоянно?};
\node[branch] (direct) at (-5.1,-11) {Прямая зависимость\\ $\downarrow\;\downarrow$\\$a_1b_2=a_2b_1$\\Решите уравнение.};
\node[branch] (other) at (0,-11) {Выберите другой метод и решите задачу.};
\node[branch] (inverse) at (5.1,-11) {Обратная зависимость\\ $\downarrow\;\uparrow$\\$a_1b_1=a_2b_2$\\Решите уравнение.};
\node[act] (check) at (0,-14) {Проверьте значение, единицы и смысл результата.};
\node[term] (end) at (0,-15.7) {Запишите ответ};
\draw[arr] (start)--(tab);\draw[arr] (tab)--(base);\draw[arr] (base)--(d);
\draw[arr] (d.west)-| node[pos=.25,above]{да} (direct.north);
\draw[arr] (d)--node[right]{нет}(e);
\draw[arr] (e.east)-| node[pos=.25,above]{да} (inverse.north);
\draw[arr] (e)--node[right]{нет}(other);
\draw[thick,draw=darkaccent] (direct.south)|-(0,-12.8)-|(inverse.south);
\draw[thick,draw=darkaccent] (other.south)--(0,-12.8);
\fill[darkaccent] (0,-12.8) circle (1.6pt);
\draw[arr] (0,-12.8)--(check.north);
\draw[arr] (check)--(end);
\end{tikzpicture}
\end{center}
\vfill
\clearpage
\section*{Тренировочный блок}
Решите задачи самостоятельно. В текстовых задачах составляйте таблицу и отмечайте направление стрелок; в конце проверяйте смысл ответа.
\ex{Базовый уровень}
\begin{enumerate}[label=\arabic*.]
\item Образуют ли отношения $3{,}6:1{,}2$ и $7{,}5:2{,}5$ пропорцию? Проверьте основное свойство пропорции.
\item Найдите неизвестный член пропорции: $2\dfrac12:1\dfrac23=3:x$.
\item Из $4{,}5$ кг семян получают $1{,}8$ кг масла. Сколько масла получится из $12{,}5$ кг семян при том же выходе?
\item Автомобиль преодолевает путь за 5 ч со скоростью 72 км/ч. За сколько часов он проедет тот же путь со скоростью 90 км/ч?
\end{enumerate}
\ex{Средний уровень}
\begin{enumerate}[label=\arabic*.,start=5]
\item Из 240 участников кружка 84 занимаются шахматами. Какой процент участников занимается шахматами?
\item В $2{,}4$ кг сплава содержится $37{,}5\%$ меди. Найдите массу меди.
\item В партии сырья содержится $48\%$ полезного вещества. Его масса равна 36 кг. Какова масса партии?
\item На выполнение плана отводилось 8 месяцев. При неизменном темпе работы какой процент этого плана можно выполнить за 11 месяцев?
\end{enumerate}
\ex{Повышенный уровень}
\begin{enumerate}[label=\arabic*.,start=9]
\item Из 320 кг свежих слив получили 88 кг чернослива. Какой процент массы свежих слив составляет масса чернослива? Какую массу свежих слив нужно взять, чтобы получить 55 кг чернослива при том же выходе?
\item На переработку отправили $80\%$ собранных оливок. Из переработанных оливок получили 72 кг масла, что составляет $60\%$ их массы. Какова масса всего урожая оливок?
\end{enumerate}
\clearpage
\section*{Ответы}
\begin{center}
\begin{tabularx}{.94\linewidth}{@{}rX@{}}
\toprule № & Ответ\\\midrule
1 & Да\\
2 & $x=2$\\
3 & 5 кг\\
4 & 4 ч\\
5 & $35\%$\\
6 & $0{,}9$ кг\\
7 & 75 кг\\
8 & $137{,}5\%$\\
9 & $27{,}5\%$; 200 кг\\
10 & 150 кг\\\bottomrule
\end{tabularx}
\end{center}
\vspace{9mm}
\begin{center}{\large\color{darkaccent} Итоговая памятка}\end{center}
\begin{tabularx}{\linewidth}{@{}>{\raggedright\arraybackslash}p{.38\linewidth}X@{}}
\toprule Сигнал в задаче & Полезное действие\\\midrule
Два отношения & Проверьте равенство перекрёстных произведений.\\
Постоянное отношение & Направьте стрелки одинаково.\\
Постоянное произведение & Направьте стрелки противоположно.\\
Процент, масса, количество & Установите, что принято за $100\%$.\\
Нужно найти целое & Разделите известную часть на долю.\\
Ответ больше $100\%$ & Проверьте, является ли процент долей целого или сравнением с базой.\\\bottomrule
\end{tabularx}
\end{document}
